\textbf{(CAP-COLUNI 2022 -- Q11)} Com o advento da pandemia da COVID-19, a população brasileira começou a ser vacinada contra o vírus, no mês de janeiro de 2021. Considere, como apta para receber a vacina, a população que tem 18 anos de idade ou mais.

Um município no interior de Minas Gerais iniciou o processo de imunização contra a COVID-19 no mês de fevereiro de 2021. Sabendo-se que a população total dessa cidade é de 22.500 habitantes. Desse total, 70\% tem idade maior ou igual a 18 anos.

A Secretaria Municipal de Saúde dessa cidade divulgou o gráfico abaixo, com a quantidade de habitantes vacinados, por mês, com a primeira dose dos imunizantes disponibilizados para o município. Sabendo-se que, até o mês de outubro de 2021, 100\% da população desse município, considerada apta, recebeu a primeira dose da vacina.

\includegraphics{./figs/2022-Q11.png}

Em decorrência de um problema, o número de vacinas referentes à primeira dose aplicadas no mês de julho foi digitado como 0 (zero), de forma errada. É CORRETO afirmar que o total de pessoas desse município vacinadas com a primeira dose, no mês de julho, equivale a:



\begin{enumerate}
	\item 12\% da população apta e 8,4\% da população total do município.
	\item 18,9\% da população apta e 4,8\% da população total do município.
	\item 8,4\% da população apta e 12\% da população total do município.
	\item 12\% da população apta e 4,8\% da população total do município.
\end{enumerate}
